%================================================================
\chapter{Introducción}
\addcontentsline{toc}{chapter}{Introducción}
\markboth{Introducción}{Introducción}
%================================================================

\begin{infobox}
Este documento \textbf{no} está pensado para leerse de principio a fin.
El profesorado y el personal que ya se sientan cómodos con \LaTeX{}, Git y
Python pueden pasar directamente a los capítulos de unidades del curso. Quienes
sean nuevos en una o más de las herramientas de apoyo deberían consultar primero
el capítulo de infraestructura pertinente. Utilice la tabla de rutas de las
páginas siguientes para identificar qué capítulos corresponden a su perfil.
\end{infobox}

%================================================================
\section*{Convenciones utilizadas en este documento}
%================================================================

\noindent
\textbf{Cuadros destacados.} En el texto aparecen cuatro tipos de cuadros destacados:

\begin{itemize}
    \item \textbf{Información} (borde verde azulado): material de referencia y
          contexto neutral.
    \item \textbf{Consejo} (borde azul): orientación práctica y
          atajos.
    \item \textbf{Advertencia} (borde magenta): precauciones críticas
          que, si se ignoran, provocarán errores o pérdida de datos.
    \item \textbf{Nota} (borde naranja): elementos que requieren atención
          pero no son críticos.
\end{itemize}

\noindent
\textbf{Listados de código.} Todos los comandos que se pretende escribir en una
terminal se muestran en bloques de código con tipografía monoespaciada. Los comandos de Windows usan
un fondo oscuro; los comandos de Linux/macOS usan un fondo más claro.
Los comandos que funcionan de forma idéntica en todas las plataformas usan el estilo
de listado neutral.

\noindent
\textbf{Referencias cruzadas.} El texto subrayado en azul es un hiperenlace
cliqueable. Las referencias a capítulos, secciones, figuras y tablas también
son cliqueables en el PDF.

%================================================================
\section*{Cómo está organizado este documento}
%================================================================

El documento se divide en dos partes. La primera parte cubre la
infraestructura técnica necesaria para trabajar con el repositorio del
currículo: configurar el entorno de software, usar la
interfaz de línea de comandos, escribir en \LaTeX{}, gestionar versiones con
Git y GitHub, y trabajar con herramientas locales de modelos de lenguaje. Estos
capítulos están escritos para profesorado y personal que quizás no tengan
formación en desarrollo de software; los colaboradores experimentados pueden ojearlos
o saltárselos por completo.

La segunda parte contiene los capítulos de unidades del curso. Cada unidad
corresponde a un curso del catálogo del departamento y documenta el
título del curso, los resultados de aprendizaje del estudiante (SLO), el esquema temático y
la estrategia de evaluación. Las unidades están numeradas según su identificador de catálogo
(por ejemplo, la unidad 1023 corresponde a \textsc{Math 1023}). Los
capítulos dentro de cada parte son independientes; un lector interesado únicamente
en un curso específico no necesita leer ningún otro capítulo de unidad.

%================================================================
\section*{Dependencias entre capítulos}
%================================================================

      \renewcommand{\thefigure}{\arabic{figure}}
      \setcounter{figure}{0}
      La Figura~\ref{fig:dependencies} muestra la estructura de dependencias de los
      capítulos de infraestructura. Una flecha continua indica una dependencia estricta:
      el capítulo de destino requiere que el de origen se haya completado primero. Una
      flecha discontinua indica únicamente un orden recomendado.

\begin{figure}[htbp]
\centering
\begin{tikzpicture}[
    node distance=1.6cm and 2.4cm,
    every node/.style={font=\small},
    req/.style={rectangle, rounded corners, draw=colorRequired,
                fill=colorRequired!15, text width=3.2cm, align=center,
                minimum height=0.9cm},
    ess/.style={rectangle, rounded corners, draw=colorEssential,
                fill=colorEssential!15, text width=3.2cm, align=center,
                minimum height=0.9cm},
    des/.style={rectangle, rounded corners, draw=colorDesirable,
                fill=colorDesirable!15, text width=3.2cm, align=center,
                minimum height=0.9cm},
    opt/.style={rectangle, rounded corners, draw=colorOptional,
                fill=colorOptional!15, text width=3.2cm, align=center,
                minimum height=0.9cm},
    hardarrow/.style={->, thick},
    softarrow/.style={->, dashed}
]
% Fila 1
\node[req] (env)  {Configuración del entorno};
\node[req, right=of env] (cli) {CLI};
\node[req, right=of cli] (prog) {Introducción a la programación};

% Fila 2
\node[ess, below=of env] (pyenv) {Entornos de Python};
\node[des, below=of cli] (numops) {Operaciones numéricas};
\node[req, below=of prog] (latex) {\LaTeX{}};

% Fila 3
\node[req, below=of pyenv] (git) {Git y GitHub};
\node[opt, below=of numops] (llm) {LLM local};

% Flechas
\draw[hardarrow] (env) -- (cli);
\draw[hardarrow] (cli) -- (prog);
\draw[hardarrow] (env) -- (pyenv);
\draw[hardarrow] (prog) -- (pyenv);
\draw[softarrow] (pyenv) -- (numops);
\draw[hardarrow] (pyenv) -- (git);
\draw[softarrow] (cli) -- (latex);
\draw[softarrow] (numops) -- (llm);
\draw[softarrow] (git) -- (llm);
\end{tikzpicture}
\caption{Grafo de dependencias para los capítulos de infraestructura. Las flechas continuas
indican dependencias estrictas; las flechas discontinuas indican el orden recomendado. El color
de los nodos codifica la categoría: \textcolor{colorRequired}{\textbf{Obligatorio}},
\textcolor{colorEssential}{\textbf{Esencial}},
\textcolor{colorDesirable}{\textbf{Deseable}},
\textcolor{colorOptional}{\textbf{Opcional}}.}
\label{fig:dependencies}
\end{figure}


\begin{landscape}
%================================================================
\section*{¿Qué capítulos necesita?}
%================================================================

Utilice la Tabla~\ref{tab:routing} para identificar su ruta de lectura según su
familiaridad con las herramientas de apoyo. Los capítulos marcados como \textbf{Obligatorio}
deben completarse antes de trabajar con el repositorio. \textbf{Ojear} significa
leer los títulos de las secciones y realizar los ejercicios, pero omitir la prosa
explicativa. \textbf{Omitir} significa que el capítulo es opcional para su perfil.

\begin{table}[htbp]
\centering
\caption{Ruta de lectura recomendada según la experiencia previa con el conjunto de herramientas.}
\label{tab:routing}
\footnotesize
\begin{tabular}{l*{8}{c}}
\toprule
\textbf{Perfil}
  & \rotatebox{70}{Config.\ entorno}
  & \rotatebox{70}{CLI}
  & \rotatebox{70}{Intro.\ prog.}
  & \rotatebox{70}{Entornos Python}
  & \rotatebox{70}{Oper.\ numéricas}
  & \rotatebox{70}{\LaTeX{}}
  & \rotatebox{70}{Git}
  & \rotatebox{70}{LLM} \\
\midrule
Nuevo en todas las herramientas
  & $\bullet$ & $\bullet$ & $\bullet$ & $\bullet$ & $\bullet$
  & $\bullet$ & $\bullet$ & O \\
Familiarizado con \LaTeX{}, nuevo en Python y Git
  & $\bullet$ & $\bullet$ & $\bullet$ & $\bullet$ & O
  & -- & $\bullet$ & O \\
Familiarizado con Python y Git, nuevo en \LaTeX{}
  & $\bullet$ & O & O & O & --
  & $\bullet$ & -- & O \\
Experimentado con todas las herramientas
  & O & -- & -- & O & --
  & -- & -- & O \\
\bottomrule
\end{tabular}

\bigskip
\noindent\footnotesize
$\bullet$ = Obligatorio \quad O = Ojear \quad -- = Omitir
\normalsize
\end{table}

%================================================================
\section*{Registro de progreso}
%================================================================

El capítulo final, Capítulo~\ref{ch:assessment} (Evaluación y registro de
progreso), organiza los elementos clave de finalización de todos los capítulos de infraestructura
en cuatro categorías: Obligatorio, Esencial, Deseable y Opcional. Abra
el formulario PDF rellenable en Adobe Acrobat Reader para registrar y guardar sus
respuestas. Utilícelo durante todo el proceso de incorporación como lista de verificación de preparación.

\newpage
%================================================================
\section*{Tiempos estimados de finalización}
%================================================================

La Tabla~\ref{tab:times} muestra los tiempos estimados de finalización por capítulo
de infraestructura. Los tiempos suponen que no se ha tenido exposición previa al tema; los colaboradores con
formación en un área determinada la completarán más rápido. Los capítulos de unidades del curso
se leen como material de referencia y no están incluidos en esta tabla; el tiempo de lectura
de una sola unidad es aproximadamente de 20 a 30 minutos.

\begin{table}[htbp]
\centering
\caption{Tiempo estimado de finalización por capítulo de infraestructura.}
\label{tab:times}
\begin{tabular}{clr}
\toprule
\textbf{Capítulo} & \textbf{Título} & \textbf{Tiempo estimado} \\
\midrule
1  & Configuración del entorno                              & 2 horas \\
2  & Interfaces de línea de comandos y scripting de shell    & 2 horas \\
3  & Introducción a los entornos de programación       & 2 horas \\
4  & Entornos virtuales de Python y automatización de flujos de trabajo & 1.5 horas \\
5  & Operaciones numéricas con NumPy                & 2 horas \\
6  & Introducción a \LaTeX{}                       & 1.5 horas \\
7  & Control de versiones con Git y GitHub            & 1.5 horas \\
8  & Infraestructura local de LLM con Ollama           & 1 hora (colaboradores) / 3 horas (administradores) \\
9  & Evaluación y registro de progreso                 & Continuo \\
\midrule
   & \textbf{Total (ruta completa)}                     & \textbf{$\approx$ 13.5--15.5 horas} \\
\bottomrule
\end{tabular}
\end{table}

\end{landscape}
